\section{Threats to Validity and Limitations}
\label{sec:limits}
% NEW (written for the TOMPECS version; not in the original paper)
We group the limitations of the study by the kind of validity they threaten.

\subsection{Measurement and Internal Validity}
\paragraph{Measurement on shared hardware.}
All experiments ran on shared GPU nodes. Identical code drifted by up to $49\%$ between sessions, so we
build every arm once, time them round-robin in a single session, wait for an idle GPU, log other users'
GPU processes at the start and end of every repetition, and report medians over clean repetitions only.
Two complete timing passes were discarded because other users' jobs occupied the GPU. Even so,
detection is imperfect: in one Cholla pass three repetitions drifted by $15\%$ while our detector
reported no foreign GPU process, so we additionally require repetitions to agree before using them.
Residual interference cannot be excluded on a shared machine: measurement noise alone puts a floor of
roughly $2\%$ under any effect we report on the applications, and far higher on the noisiest benchmarks,
where per-repetition spread exceeds $100\%$.

\paragraph{Tooling defects discovered during the study.}
The prompt-fidelity defects of Section~\ref{sec:defects} (Table~\ref{tab:defects}) were all present
during the benchmark study: profiling summaries reached the model with wrong units, truncated
stall data and understated launch counts, and the repair channel discarded the diagnostics a model
needed. They are fixed, but every result predating a given fix was produced with the affected prompt
and would need regeneration to be strictly comparable. We report the benchmark numbers as measured
rather than re-running the grid, and we do not claim that the ablation of Section~\ref{sec:ablation}
would reproduce unchanged on the corrected summarizers.

\subsection{Statistical Conclusion Validity}
\paragraph{Sampling of model outputs.}
Except where stated, each (application, model, prompt) cell is a \emph{single} LLM sample. Language
model outputs vary between calls, so a single cell conflates prompt effect with sampling variance.
% EDIT (2026-10-26): the earlier text claimed a stable replicate of the structural/diagnostic arms.
% That round (report sec. 12) was timed on the H100 node while its driver state produced wrong
% output and 20-300x inflated times (report sec. 14), covered 6 GPT-5 pairs before it was stopped,
% and used a revised checklist; it is not evidence either way.
Only one setting is replicated: a second, independently generated sample of the baseline prompt on
both applications, which keeps Cholla on the same side of the profile-data comparison
(Section~\ref{sec:profabl}) and reproduces AutoDock-GPU's tie. The prescriptive-guidance comparison of
Section~\ref{sec:guidvar} has no valid replicate: a second generation round was timed on the H100 node
while a driver change there made unmodified benchmarks produce wrong output and run $20$--$300\times$
slower, and it was stopped after six pairs. Its cells, like the per-cell numbers in
Section~\ref{sec:results}, are single samples with no confidence interval. Combined
with the measurement floor above, this means an effect below roughly $5\%$ in a single cell should be
treated as unresolved rather than small.

\paragraph{Coverage of the profile-data ablation.}
% EDIT (2026-10-27): counts from the CUDA 13.3.1 re-timing; boxfilter and ccsd-trpdrv now timed.
The comparison in Section~\ref{sec:profabl} spans 24 (application, model) pairs and scores 23:
\texttt{background-subtract} under GPT-5 yields no correct \texttt{code\_only} candidate. The positive
result rests on a single application, Cholla, with two models; the nineteen benchmark pairs show that
profile data moves individual results in both directions there with no net gain, not that the gain on
Cholla generalizes to other large applications.

\subsection{Construct Validity}
\paragraph{Correctness gates are empirical, not proofs.}
Each application is gated by its own output: bit-exact field comparison where the application is
deterministic (Cholla) and a statistical band where it is not (AutoDock-GPU, whose GPU results
vary run to run even with a fixed seed). We validated every gate by injecting faults, and the Cholla
gate detects a $10^{-6}$ relative perturbation of a single flux. Nevertheless a subtle accuracy loss
that stays inside the AutoDock band would pass, gates check the outputs the application prints rather
than all internal state. We gate every build that is timed, but ``gated correct'' means
``passed these checks'', not ``verified''.

\paragraph{Benchmark gaming is detected heuristically.}
Transformations that delete the measured work while preserving printed output defeat output comparison
by construction (Section~\ref{sec:intent}). We flag speedups above $5\times$ and inspect the diff. This
catches the large cases we observed ($1504\times$, $2809\times$, $235$--$949\times$, $31\times$), but a
gaming transformation yielding, say, $1.5\times$ would not be flagged automatically, and we cannot rule
out such cases in the reported numbers.

\paragraph{A node fault invalidated the first benchmark timing.}
% EDIT (2026-10-27): replaces "Benchmark verification was unusable as a gate on that node". The FAIL
% verdicts it described were caused by the driver change (real-applications report, sec. 14); on the
% CUDA 13.3.1 re-timing every unmodified benchmark with a check prints PASS.
The HeCBench half of the profile-data ablation was first timed after a driver update on the shared
H100 node, which made binaries built with CUDA 12.x print wrong output and run $20$--$300\times$
slower without failing. Unmodified benchmarks then printed FAIL from their own checks, and two
exceeded our $600$\,s per-run cap. We detected this only by comparing against earlier measurements of
the same binaries, discarded that timing, and re-timed the same candidates rebuilt with CUDA 13.3.1;
Section~\ref{sec:profabl} reports only the re-timing. The candidates themselves were generated on the
faulty node, so a candidate rejected by the generation-time check could have been rejected wrongly; the
one case where this decides a result (\texttt{background-subtract}, GPT-5) we re-checked with the
working toolchain, and the rejection stands. A benchmark's built-in check is also not always an oracle:
\texttt{chemv}'s timing loop launches its kernels $1000$ times, accumulating into the same output
buffer, and then compares against a single application of the CPU reference, so the check cannot pass
for the unmodified program; we therefore exclude it (Section~\ref{sec:guidvar}). For the three
benchmarks that print no verdict (\texttt{amgmk}, \texttt{atomicSystemWide}, \texttt{bincount}),
correctness rests on the pipeline's generation-time check and on our inspection of the diffs, which is
weaker than the gates we built for the applications.

\paragraph{Scope of the objective.}
We optimize execution time of a single GPU kernel file under a correctness gate. We do not study
energy, memory footprint, multi-GPU scaling, numerical stability under long runs, or maintainability of
the generated code---all of which matter for production use of these applications.

\subsection{External Validity}
\paragraph{Hardware and software scope.}
Results cover NVIDIA GPUs only: H100 PCIe (sm\_90) for the benchmark studies and GH200 (sm\_90) for
the applications, with CUDA 12.6.3 (13.3.1 for the re-timed half of the profile-data ablation).
Conclusions may not transfer across architectures, and we did not test that
directly: every prompt variant was evaluated on one architecture per study, so an effect that depends
on register file size, SM count or cache hierarchy would be invisible to us. The profiles supplied for
the six original ablation benchmarks were themselves collected on an A100 and used on an H100. We use two proprietary models whose weights and versions we do not control; the
same prompt may behave differently after a model update.

\paragraph{Single-file editing.}
The pipeline gives the model one source file and replaces that file. The model therefore cannot change
headers, launch sites or build flags, and---more importantly---cannot see how its file is compiled.
This caused three of our real-application failures (symbol redefinition in a textually included file,
and kernel-signature changes that broke linkage against ten call sites). Some optimizations a human
would make are out of reach by construction.
